\documentclass[11pt]{article}
\usepackage[T1]{fontenc}
\usepackage[utf8]{inputenc}
\usepackage{lmodern}
\usepackage[margin=1in]{geometry}
\usepackage{amsmath,amssymb,amsthm,mathtools}
\usepackage{booktabs,tabularx,array}
\usepackage{enumitem}
\usepackage{microtype}
\usepackage[colorlinks=true,linkcolor=blue,citecolor=blue,urlcolor=blue]{hyperref}
\usepackage{tikz-cd}
\newtheorem{theorem}{Theorem}[section]
\newtheorem{proposition}[theorem]{Proposition}
\newtheorem{lemma}[theorem]{Lemma}
\newtheorem{definition}[theorem]{Definition}
\newtheorem{remark}[theorem]{Remark}
\DeclareMathOperator{\im}{im}
\DeclareMathOperator{\conv}{conv}
\DeclareMathOperator{\spanlin}{span}
\DeclareMathOperator{\rank}{rank}
\DeclareMathOperator{\nnrank}{rank_+}
\DeclareMathOperator{\Alg}{Alg}
\DeclareMathOperator{\Fact}{Fact}
\newcommand{\Set}{\mathbf{Set}}
\newcommand{\Conv}{\mathbf{Conv}}
\newcommand{\FinStoch}{\mathbf{FinStoch}}
\newcommand{\RR}{\mathbb R}
\newcommand{\HH}{\mathcal H}
\newcommand{\DD}{\mathsf D}
\newcommand{\VV}{\mathsf V}
\newcommand{\PP}{\mathsf P}
\setlist{nosep,leftmargin=*}
\setlength{\parskip}{3pt}
\title{A common generated observable construction and its native boundaries}
\author{Technical audit of the categorical phenomenology programme}
\date{2 November 2026}
\begin{document}
\maketitle
\begin{abstract}
A genuine common construction already underlies several of the programme's examples: extend preparation profiles freely in a supplied algebraic doctrine, then take their image in a supplied response algebra. It gives the ordinary response quotient, meet completion, linear span and convex hull, with typed comparisons under compatible changes of tests and doctrine. This is a substantive common fragment, not merely similar terminology. It is also already present in the recovered notes and established categorical mathematics. The wider material is best described as a structured family of mechanisms: predictive filtering, positive realization, constrained latent-state minimization and cylinder normalization are not all the same universal problem.

The audit makes two boundaries explicit. Test restriction may preserve every distinction between observed preparations while losing distinctions introduced by completion. In finite stochastic realization, the generated convex image is representable by a finite simplex with both algebraic reachability and observability exactly when it is a simplex. This does not prohibit latent realizations of nonsimplicial images. Two affinely isomorphic square images have minimum latent sizes four and three in their respective output embeddings. The conclusions consolidate and check existing work; no novelty claim or solution of the full programme is made.
\end{abstract}

\section{Scope and the recovered baseline}

The question here is whether there is a precise shared construction worth developing, rather than whether all observationally specified mathematics can be declared an instance of one unrestricted theory. The starting data are preparations, observations and response profiles. Algebraic operations, admissible realizations and comparison maps must be separately supplied. Calling a convex combination or a formal meet a state does not establish its empirical availability.

The source recovery materially changes the assessment. The existing Boolean discussion already gives the free-image theorem, its response-quotient/FCA comparison, the extension criterion in terms of relations between generators, and a concrete obstruction under changed tests. The stochastic discussion already distinguishes response quotients, convex images and minimum latent factorizations, including a rank/nonnegative-rank gap and nonunique minimum factorizations. The positive-Hankel discussion already treats invariant cones and exact comparison spans. The completed cylinder and cross-doctrine checkpoints already prove coherent normalization spans and local exactness criteria. These are not missing results to be rediscovered.

\begin{center}
\small
\begin{tabularx}{\textwidth}{@{}p{0.26\textwidth}Xp{0.24\textwidth}@{}}
\toprule
Claim or mechanism & Recovered status & Status in this report\\
\midrule
Free extension followed by image & Explicit in the Boolean discussion and problem note & Reproved and typed uniformly\\
Response quotient and FCA comparison & Explicit, including extension obstruction & Recovered; small stress case checked\\
Convex image versus stochastic latent minimum & Explicit, including rank gap & Recovered; simplex boundary isolated\\
Positive Hankel invariant cones & Explicit, using classical positive realization theory & Retained as a separate mechanism\\
Cylinder and cross-doctrine spans & Completed checkpoints, not an open generic target & Credited; not replaced here\\
Paired square embeddings and mixed square & Concrete audit examples & Verified illustrations, not novelty claims\\
\bottomrule
\end{tabularx}
\end{center}

The recovered source ledger distinguishes uploaded exploratory conversations from completed checkpoints. This report proves the elementary claims it uses below, but does not certify every theorem asserted in those conversations. In particular, no general positive-cone closedness, tensor-rank bound or physical interpretation is imported without its own hypotheses. The primary-source review in Section~\ref{sec:prior} is scoped, not an exhaustive novelty search.

\section{One construction with four concrete instances}
\label{sec:construction}

Let $M$ be a monad on $\Set$ that preserves surjections. Write its unit and multiplication as $\eta$ and $\mu$. Let $(B,b:MB\to B)$ be an $M$-algebra, let $P$ be a set of preparations, and let
\[
 r:P\longrightarrow B
\]
be the observed profile map. Typically $B$ is a specified subalgebra of a product $\Omega^E$, where $E$ is a set of tests. Neither $M$ nor the algebra structure on $B$ is inferred from the table $r$.

\begin{definition}
The \emph{generated observable image} is
\[
 r^\#=b\,Mr:MP\longrightarrow B,
 \qquad \HH_M(r)=\im(r^\#)\subseteq B.
\]
Write $r^\#=mq$, where $q$ is surjective and $m$ is inclusion. The distinguished preparations in $\HH_M(r)$ are $q\eta_P$.
\end{definition}

\begin{theorem}[Recovered generated-image mechanism]
The set $\HH_M(r)$ has a unique $M$-algebra structure making $m$ a homomorphism. It is the smallest $M$-subalgebra of $B$ containing $r(P)$, and $q$ is a homomorphism.

A realization of $r$ consists of an $M$-algebra $A$, a map $a:P\to A$, and a homomorphism $o:A\to B$ with $oa=r$. A realization morphism $f:(A,a,o)\to(A',a',o')$ is a homomorphism satisfying $fa=a'$ and $o'f=o$. Call a realization \emph{generated} if $a^\#:MP\to A$ is surjective, and \emph{observable} if $o$ is injective. Then $\HH_M(r)$ is initial among observable realizations and terminal among generated realizations. It is the unique generated observable realization up to unique realization isomorphism.
\end{theorem}

\begin{proof}
Because $Mq$ is surjective and $r^\#$ is a homomorphism,
\[
 b\,Mm\,Mq=b\,M(r^\#)=r^\#\mu_P=mq\mu_P.
\]
Thus $b\,Mm$ takes values in $\HH_M(r)$ and corestricts to an algebra structure there. Composing its unit and multiplication identities with the injective $m$ verifies the algebra laws. The same calculation gives the homomorphism law for $q$. A subalgebra containing $r(P)$ contains every value of $r^\#$, proving minimality.

If $o$ is injective, identify $A$ with its subalgebra image in $B$. This image contains $r(P)$ and therefore $\HH_M(r)$. The inclusion, transported back along $o$, is the unique realization map $\HH_M(r)\to A$. If $a^\#$ is surjective, then
\[
 o(A)=\im(oa^\#)=\im(r^\#)=\HH_M(r).
\]
Corestricting $o$ gives the unique realization map $A\to\HH_M(r)$. Uniqueness in both cases follows from the indicated observation inclusions. When both conditions hold, the resulting maps are inverse.
\end{proof}

These statements concern specified subcategories of realizations. They are not an assertion that every intuitive notion of a smallest implementation is an image, or that the resulting object lies in a separately constrained implementation class.

\subsection{Ordinary response profiles}
For $M=\mathrm{Id}$, the image is simply
\[
 Q_r=r(P)\cong P/{\sim},\qquad p\sim p'\ \Longleftrightarrow\ r(p)=r(p').
\]
No new operations are introduced. If the tests are merely a fixed collection, this is a response quotient, not automatically an automaton or a predictive state space. Continuation operations must be supplied and shown to descend before such interpretations are available.

\subsection{Meet completion and formal concepts}
Let $B=\{0,1\}^E$, identified with subsets of $E$, and use the finite-powerset monad interpreted by finite meets, including the empty meet $\top$. Then
\[
 L_r=\left\{\bigcap_{p\in U}r(p):U\subseteq P\text{ finite}\right\}.
\]
For finite $P$, this is the family of all intersections of row intents. It is the intent presentation of the FCA concept lattice; with the inclusion order used here, it is opposite to the usual extent-ordered concept lattice. For infinite $P$, finite meets do not suffice for arbitrary intersections: use the full powerset monad and arbitrary-meet algebras for that claim.

The map $Q_r\hookrightarrow L_r$ is an inclusion of observed profiles into generated profiles. It is an isomorphism exactly when the rows are already meet-closed, including the empty meet. The new profiles are consequences of the chosen meet doctrine, not newly measured preparations.

\subsection{Linear and convex completion}
For the free real-vector-space monad $\VV$ and $B=\RR^E$,
\[
 \HH_{\VV}(r)=\spanlin\{r(p):p\in P\}.
\]
For the finitely supported probability-distribution monad $\DD$ and a convex response algebra $B$,
\[
 \HH_{\DD}(r)=\conv\{r(p):p\in P\}.
\]
A simplex or a product of simplices can be such a response algebra. In the convex case, the ambient category is the category of convex algebras and affine maps, not just the finite free algebras $\Delta S$.

The monad morphism $\DD\Rightarrow\VV$ interprets a probability distribution as a formal real linear combination. With a common real response space it gives the affine inclusion
\[
 \conv(r(P))\hookrightarrow\spanlin(r(P)).
\]
It does not make signed or unnormalized combinations operationally executable. Nor is there a canonical translation between every pair of the four doctrines. In particular, shared profile notation does not supply a meet-to-linear or linear-to-convex monad morphism.

Even empty preparation sets reveal the supplied structure: the ordinary image is empty, the linear image is $\{0\}$, and the meet image contains $\top$. The convex image is empty if empty convex algebras are admitted. All native stochastic statements below assume finite nonempty input and output sets.

\section{Changing doctrine and observations with correct types}
\label{sec:changes}

Let $S,M$ be monads as above, and let $\alpha:S\Rightarrow M$ be a monad morphism. Its reduct functor is
\[
 \alpha^*:\Alg(M)\longrightarrow\Alg(S),
 \qquad (C,c)\longmapsto(C,c\alpha_C).
\]
The direction matters: an $M$-algebra is regarded as an $S$-algebra, not conversely.

\begin{proposition}[Concrete comparison data]
Suppose $r:P\to(B,b)$ is an $S$-presentation and $s:Q\to(C,c)$ an $M$-presentation. Given a preparation map $g:P\to Q$ and an $S$-homomorphism $h:B\to\alpha^*C$ with $hr=sg$, there is a unique induced $S$-homomorphism
\[
 k:\HH_S(r)\longrightarrow\alpha^*\HH_M(s)
\]
whose composite with the target profile inclusion is $h$ restricted to $\HH_S(r)$. The comparisons preserve identities and compose coherently.
\end{proposition}

\begin{proof}
The homomorphism law, naturality of $\alpha$, and $hr=sg$ give
\[
 h r_S^\#=s_M^\#\,Mg\,\alpha_P.
\]
Hence $h$ sends the source image into the target image, and its restriction is an $S$-homomorphism. The target inclusion gives uniqueness. For consecutive comparisons, both the constructed comparison and the composite are the restriction of the composite response map. This proves coherence, with the usual reduct identifications.
\end{proof}

This is a sufficient input-level specification of a change, not a claim that every change of scientific vocabulary is already so specified. The target is explicitly a reduct. There is no fictitious common algebra category in which different structures silently become identical.

\subsection{Test restriction is surjective but need not be exact}
Fix a response algebra $\Omega$. If $E\subseteq E'$ are retained and fine tests, precomposition gives a homomorphism
\[
 \rho:\Omega^{E'}\longrightarrow\Omega^E.
\]
For $r_E=\rho r_{E'}$, it induces a surjection
\[
 q_M:\HH_M(r_{E'})\twoheadrightarrow\HH_M(r_E),
\]
since every generated coarse profile is the restriction of the same formal combination of fine rows. Thus restriction and reconstruction have a canonical comparison without preserving all distinctions.

For $\alpha:S\Rightarrow M$, the unit-like doctrine comparisons give inclusions $j$ and a commuting square, with the right-hand objects understood as $S$-reducts:
\[
\begin{tikzcd}[column sep=large]
\HH_S(r_{E'}) \arrow[r,hook,"j_{E'}"] \arrow[d,two heads,"q_S"'] &
\HH_M(r_{E'}) \arrow[d,two heads,"q_M"]\\
\HH_S(r_E) \arrow[r,hook,"j_E"'] & \HH_M(r_E).
\end{tikzcd}
\]
Here the notation presupposes that the $S$-response structure is the $\alpha$-reduct of the $M$-structure. In the concrete subset presentation this square is a pullback exactly when
\begin{equation}
 q_M^{-1}\bigl(\HH_S(r_E)\bigr)=\HH_S(r_{E'}).
 \label{eq:pullback}
\end{equation}
Indeed, its set-theoretic pullback consists precisely of fine $M$-profiles whose coarse restriction is $S$-generated. Since the forgetful functor from $S$-algebras creates limits, this also tests the pullback in $\Alg(S)$.

\paragraph{A two-row obstruction.}
Take $P=\{a,b\}$, fine tests $E'=\{x,y\}$, and rows $r(a)=10$, $r(b)=01$. Let $S=\mathrm{Id}$ and $M$ be meet completion with top. Then
\[
 Q_{E'}=\{10,01\},\qquad L_{E'}=\{00,10,01,11\}.
\]
Retaining only $x$ gives $Q_E=L_E=\{0,1\}$. Both $q_S$ and the bottom inclusion are bijective: the retained test still distinguishes the original preparations. Nevertheless the pullback contains all four fine meet profiles, while the upper left contains only two. The profiles $00$ and $11$ imitate old rows on the retained test. Coherence holds; cartesianness fails. A test adequate for the observed rows need not be adequate for structure introduced by completion.

\subsection{A checkable sufficient exactness certificate}
\begin{proposition}
Suppose that for each omitted test $e'\in E'\setminus E$ a homomorphism $f_{e'}:\Omega^E\to\Omega$ is supplied and satisfies
\[
 r_{E'}(p)(e')=f_{e'}(r_E(p))\qquad(p\in P).
\]
Then $q_M$ is an isomorphism. The same holds in every reduct doctrine $S\Rightarrow M$, and the mixed square is a pullback.
\end{proposition}
\begin{proof}
Together with the retained coordinate projections, the $f_{e'}$ define a homomorphism $F:\Omega^E\to\Omega^{E'}$ satisfying $\rho F=1$ and $Fr_E=r_{E'}$. The maps $F\rho$ and the identity agree on fine generators, hence agree on the generated fine algebra. Restrictions of $F$ and $\rho$ are inverse, also after reduct.
\end{proof}

For linear completion, omitted columns that are linear combinations of retained columns supply such a certificate. For convex completion, affine expressions do so, provided their codomain is the actual admissible response object. For Boolean meet completion, conjunctions of retained coordinates, including the empty conjunction, preserve meets and top. In the two-row obstruction, the omitted $y$ is Boolean negation of $x$ on the observed rows. Negation is a set map, but not a meet/top homomorphism. That is exactly why the same retained column is sufficient in one doctrine and insufficient in the other.

This certificate is sufficient, not generally necessary: an inverse defined on the generated subalgebra need not extend to the whole response algebra. Verifying finitely many row equations does not by itself verify that an arbitrary proposed function is a homomorphism. In concrete linear, affine or finite Boolean descriptions the latter is a separate tractable check; no general decidability or noisy-data inference theorem is claimed.

\section{Finite stochastic realization and the native class}
\label{sec:stochastic}

Let $H$ be a row-stochastic $|P|\times |Y|$ matrix, with finite nonempty $P,Y$. It defines an affine map $\bar H:\Delta P\to\Delta Y$. Its generated observable convex realization is
\[
 C_H=\conv\{H_p:p\in P\}\subseteq\Delta Y.
\]
A finite latent stochastic realization is a factorization $H=AB$, where $A:P\to\Delta S$ and the rows of $B$ belong to $\Delta Y$. Its state algebra is the finite free convex algebra $\Delta S$. The latent emission rows need not belong to $C_H$: one only needs
\[
 C_H\subseteq\conv\{B_s:s\in S\}.
\]
This is why a generated image and a smallest latent implementation are different questions.

\begin{theorem}[Exact finite-simplex boundary]
There is a finite stochastic realization $H=AB$ for which $\bar A:\Delta P\to\Delta S$ is surjective and $\bar B:\Delta S\to\Delta Y$ is injective if and only if $C_H$ is a simplex. Equivalently, the distinct extreme rows of $H$ are affinely independent.
\end{theorem}
\begin{proof}
If both conditions hold, then $\bar B(\Delta S)=\bar B\bar A(\Delta P)=C_H$ and $\bar B$ is an affine isomorphism $\Delta S\cong C_H$.

Conversely, let $S$ index the vertices of the simplex $C_H$, and let $B_s$ be those vertices. Every vertex of a finite row hull is a row of $H$. Express each row in its unique barycentric coordinates to define $A$. Each vertex row gives a preparation mapped to the corresponding $\delta_s$, so $\bar A$ is surjective. Affine independence of the vertices makes $\bar B$ injective.
\end{proof}

Here ``generated'' means surjective on the entire convex algebra, not merely that every latent state has positive probability somewhere. It also does not mean ``epimorphism in $\FinStoch$'': that is a different condition. With row-stochastic matrices, categorical epis have full column rank and monos have full row rank. A full-column-rank preparation matrix can have a proper convex image in its target simplex. For clarity, the rank characterizations follow by cancellation when the relevant rank is full. If a stochastic matrix $F$ has a nonzero right kernel vector $v$, perturb the two columns of a strictly positive binary channel by $+\epsilon v$ and $-\epsilon v$; for sufficiently small $\epsilon$, distinct stochastic channels have the same composite with $F$, so $F$ is not epic. If a nonzero row vector $w$ satisfies $wF=0$, then $w\mathbf1=wF\mathbf1=0$, so opposite small perturbations of a strictly positive input distribution by $w$ are distinct probability distributions with the same image, and $F$ is not monic.

\subsection{A canonical comparison that stays ambient}
For any factorization $H=AB$, let $R_A=\conv\{A_p:p\in P\}$. There is a canonical span in $\Conv$,
\begin{equation}
 \Delta S\ \longleftarrow\ R_A\ \longrightarrow\ C_H,
 \label{eq:span}
\end{equation}
with inclusion on the left and the surjective restriction of $\bar B$ on the right. It first restricts to algebraically reachable latent distributions, then identifies distributions with the same response.

This span is natural, not a choice of bases or a selection of a minimum factorization. A stochastic factorization morphism $F:S\to\Delta S'$ satisfying
\[
 AF=A',\qquad FB'=B
\]
in matrix notation sends $R_A$ onto $R_{A'}$, because affine maps preserve convex hulls. Its restriction commutes with both legs of~\eqref{eq:span}; restrictions preserve identities and composition. For an output garbling $G:\Delta Y\to\Delta Z$, the preparation hull is unchanged and the right endpoint becomes $C_{HG}=G(C_H)$. The induced surjection is $G|_{C_H}$, and successive garblings compose coherently.

Neither $R_A$ nor $C_H$ need be a finite free simplex. Thus this is a genuine bridge from native factorizations to an ambient generated realization, not a normalization internal to $\FinStoch$. It also does not select a minimum latent cardinality. The same restriction/corestriction argument works for the generated subalgebras in Section~\ref{sec:construction}.

\subsection{Nonnegative rank answers the cardinality question}
For one normalized response experiment, minimum latent cardinality equals $\nnrank(H)$. To see the normalization step, suppose $H=UV$ with $U,V\geq0$. Remove zero rows of $V$ and their corresponding columns of $U$, and let $d_s=\sum_yV_{sy}>0$. Set
\[
 B_{sy}=V_{sy}/d_s,\qquad A_{ps}=U_{ps}d_s.
\]
Then $AB=H$, every row of $B$ sums to one, and every row of $A$ sums to one because the rows of $H$ do. The converse is immediate. A flattened table of several separately normalized experiments requires extra normalization constraints; the unqualified one-matrix statement should not be transferred to that setting.

Even minimum cardinality does not identify a unique realization. For
\[
 K=\begin{pmatrix}4/5&1/5\\1/5&4/5\end{pmatrix},
\]
 both $(I,K)$ and $(K,I)$ are minimum two-state factorizations. They are not isomorphic in the factorization category: reversing the comparison requires $K^{-1}$, which has negative entries. The category still has the trivial domain factorization as an initial object and codomain factorization as a terminal object. The point is not nonexistence of universal objects there, but that minimum cardinality alone does not characterize a unique generated observable object.

\section{Two square images with different latent minima}
\label{sec:squares}

Consider the already familiar boundary example, up to row and column permutations,
\[
 H_{\partial}=\frac12
 \begin{pmatrix}
 1&1&0&0\\0&1&1&0\\0&0&1&1\\1&0&0&1
 \end{pmatrix}.
\]
Its rows are four vertices of a square: opposite sums agree, while two adjacent differences are linearly independent. Thus its affine dimension is two and $\rank(H_{\partial})=3$. The generated hull is not a simplex.

\begin{proposition}
$\nnrank(H_{\partial})=4$. Hence this matrix has no epi--mono factorization internal to $\FinStoch$, and no finite-simplex generated observable realization, although it has a four-state latent realization.
\end{proposition}
\begin{proof}
The four positive diagonal positions form a fooling set: for each pair, at least one crossed entry is zero. The positive support of a nonnegative rank-one matrix is a rectangle. In a nonnegative decomposition of $H_{\partial}$, each summand's support lies in the support of $H_{\partial}$, so one summand cannot cover two of those diagonal positions. At least four summands are necessary; $H_{\partial}=I_4H_{\partial}$ gives the matching upper bound.

An epi--mono factorization through $k$ finite stochastic states would have both factors of rank $k$, hence $\rank(H_{\partial})=k$. It would force $k=3$, contradicting the nonnegative-rank bound. The generated observable claim also follows directly from the simplex criterion.
\end{proof}

Now compare
\[
 H_{\mathrm{in}}=\frac1{12}
 \begin{pmatrix}
 5&4&3\\5&2&5\\3&6&3\\3&4&5
 \end{pmatrix}.
\]
The rows form a parallelogram of affine dimension two, hence an abstract convex set affinely isomorphic to the preceding square. They are not asserted to form a Euclidean square. If its first three rows are $a,b,c$, the fourth is $b+c-a$. All four are extreme. The ordinary rank is three, and
\[
 H_{\mathrm{in}}=H_{\mathrm{in}}I_3
 \quad\Longrightarrow\quad
 3=\rank(H_{\mathrm{in}})\leq\nnrank(H_{\mathrm{in}})\leq3.
\]
Thus a nonsimplex generated image can have a rank-sized native latent realization. The three latent emissions here are the vertices of $\Delta^2$, all outside the generated row hull. It is observable but not algebraically generated. Conversely, the displayed realization $(I_4,H_{\partial})$ is generated but not observable: the two equal mixtures of opposite vertices give the same output.

For completeness, an affine isomorphism from the first square to the second, in zero-based coordinates $x_0,x_1,x_2,x_3$, is
\[
 x\longmapsto a+2x_2(b-a)+2x_3(c-a).
\]
It sends the boundary rows in order to $a,b,b+c-a,c$, and its two independent direction vectors make it an affine isomorphism on their affine hulls. The exact-arithmetic checks verify this and its inverse on spanning vertices.

\paragraph{What this does and does not show.}
These are different response tables in different output simplices. The abstract convex carrier alone does not determine minimum latent size; its embedding in the response simplex matters. The full generated realization $(C_H,r,\iota)$ retains that embedding and the entire response table. The examples do \emph{not} show that a richer common construction could never retain the information needed for latent optimization. They show that the universal image property and the optimization problem are distinct, and that shape alone is insufficient.

Changing observations can also change native membership. Grouping outputs of $H_{\partial}$ as $\{1,2\}\mid\{3,4\}$ sends its four rows to binary probabilities $1,1/2,0,1/2$. The new hull is a segment, hence a simplex, although the fine hull was not. Exact checks of all sixteen deterministic four-to-two output maps give four point hulls and twelve segment hulls, with corresponding one- or two-state generated observable factorizations. This is exhaustive for those sixteen maps, not for all stochastic garblings or dynamical observation changes.

\section{Where the wider mechanisms agree and differ}
\label{sec:prior}

The common generated-image fragment is genuine: the same formula and proof apply to four concrete constructions, and compatible test/doctrine changes yield explicitly typed maps. That warrants more than an umbrella label. But it does not license flattening all recovered mechanisms into this fragment.

\begin{enumerate}
\item \textbf{Predictive filtering has a different universal direction.} The recovered transducer discussion treats filtering as a right adjoint under specified intertwining maps, followed by reachable restriction and observable minimization. Hiding an output can require filtering again and can turn a finite predictor into an infinite one. A static restriction map on profile algebras is not that theorem.
\item \textbf{Positive linear realization adds cone admissibility and invariance.} The recovered Hankel discussion already gives generated and test-positive cone bounds and invariant-polyhedral or simplicial criteria for finite or rank-sized positive realization. Static convex-hull geometry neither replaces those dynamic invariant-cone conditions nor establishes that all cones under discussion are closed.
\item \textbf{Native latent minimization is a constrained lifting problem.} A surrounding simplex may contain the generated hull strictly. Restriction followed by an observable quotient gives an ambient comparison, but can leave the native free class. The simplex theorem locates one exact boundary; it does not solve arbitrary nonnegative factorization or dynamic positive realization.
\item \textbf{Cylinder normalization is broader than the monadic image presentation.} The recovered completed checkpoints already construct native/free normalization spans, their substitution coherence and cross-doctrine comparisons under precise factorization and mixed-pullback hypotheses. They also distinguish local exactness from global class preservation. They should not be reduced to the claim that one can draw a factorization diagram.
\end{enumerate}

The external literature is likewise stronger than a bare unifying analogy. Garner's Isbell-monad theorem identifies cylinder-factorization categories with Isbell pseudoalgebras; its lax and colax extensions already organize comparisons that preserve one factorization class, and its restricted doctrines vary the completion~\cite{Garner}. Stacey's comparison of test-defined smooth structures varies tests, underlying categories and forcing conditions, with concrete comparison criteria~\cite{Stacey}. These are substantive precedents for changing observational vocabulary and requirements.

Shen and Zhang give functorial concept-lattice constructions from contexts and infomorphisms, factored through closure spaces~\cite{ShenZhang}. Melli\`es and Zeilberger distinguish direct preservation from preservation after a specified dualization/completion in refinement systems~\cite{MelliesZeilberger}. CALF and algebraic automata learning already treat acquisition, finite observation tables and monad-parametric algebraic structure~\cite{CALF,UrbatSchroeder}. Cohen and Rothblum supply the established nonnegative-rank framework for the static native optimization boundary~\cite{CohenRothblum}. Representable Markov categories provide an established treatment of statistical informativeness and garbling comparisons~\cite{FritzEtAl}.

None of these citations automatically certifies every operational interpretation. Conversely, the existence of a parameterized doctrine, a noninvertible coherent comparison, or an image-style universal property is not by itself a new contribution. The bounded audit finds no basis for claiming novelty of that generic machinery.

\section{Verdict and implications for a later research contract}

\textbf{Verdict: a genuine common fragment inside a structured family.} For the ordinary, meet, linear and convex examples, ``free extension followed by generated observable image'' is literally one construction, not a metaphor. Its comparison laws identify exactly which supplied operations and tests are being respected. Across the full recovered programme, the appropriate description is a structured family with real bridges, not yet one demonstrated universal native realization theory. It is also not merely an umbrella: there are proofs, coherent comparisons and explicit failure conditions.

The strongest correction is historical. Much of the proposed architecture already exists in the project notes, and strong general machinery exists in the literature. Reproving it can clarify the interface, but cannot be presented as filling an unoccupied gap. The present work is an audited consolidation and a sharpened boundary statement, with independently reproduced finite examples.

Three distinctions should control subsequent interpretation. First, observations determine rows; the doctrine supplies the operations used to generate additional structure. Second, a change of tests must preserve distinctions in the generated structure, not merely distinguish the originally observed preparations, if exact transport is required. Third, an ambient universal realization may leave the native realization class, and minimizing native resources is an additional question.

No larger contract is adopted here. If further work is authorized, it should name a specific native class, supplied observations and doctrine, the desired downstream operation or optimization, and the exact comparison claimed. A useful deliverable would then be a proof or counterexample that improves on the already recovered generated-image, invariant-cone and cylinder-comparison results. The present audit does not establish that a general positive native bridge is impossible or that an alleged missing bridge is absent from the literature. It does establish that ``derive a common construction'' and ``derive coherent change of doctrine'' are too coarse to describe the remaining work honestly.

\appendix
\section{Proof and reproduction status}

The main generated-image, comparison, exactness-certificate, simplex-boundary, natural-span and matrix arguments are proved above. The supplied standard-library Python checker uses exact rational arithmetic. It verifies the Boolean profile sets; ranks, affine dimensions and exposed vertices of both matrices; the fooling-set premises; inverse affine maps on spanning vertices; the non-pullback witnesses; and all sixteen binary output maps and sixty-four two-stage deterministic garbling checks. A clean-copy run and a run with Python optimization enabled both match the expected JSON output. These are finite certificates supporting the displayed elementary arguments, not a formal proof of the categorical theorems, an exhaustive novelty review or a test of physical realizability.

The accompanying source ledger records the recovered project evidence, including the Boolean/stochastic conversation, the positive-Hankel conversation, the cylinder/wiring checkpoint, the cross-doctrine checkpoint and the original problem note. The named finite checks add reproducibility without changing the attribution of inherited results. The report is intentionally delivered as uncompiled \LaTeX\ source; pagination and rendered layout have not been certified.

\begin{thebibliography}{99}
\bibitem{Garner} R. Garner, \emph{The Isbell monad}. Theorems 4.1 and 5.1; restricted doctrines in Section 6. \url{https://arxiv.org/abs/1410.7108}.
\bibitem{Stacey} A. Stacey, \emph{Comparative Smootheology}. Section 4, including Proposition 4.17. \url{https://arxiv.org/abs/0802.2225}; published version: \url{https://tac.mta.ca/tac/volumes/25/4/25-04abs.html}.
\bibitem{ShenZhang} L. Shen and D. Zhang, \emph{The concept lattice functors}. \url{https://doi.org/10.1016/j.ijar.2012.07.002}; author text: \url{https://math.mickeylili.com/files/preprints/The_concept_lattice_functors.pdf}.
\bibitem{MelliesZeilberger} P.-A. Melli\`es and N. Zeilberger, \emph{An Isbell Duality Theorem for Type Refinement Systems}. \url{https://noamz.org/papers/isbell-final.pdf}.
\bibitem{CALF} G. van Heerdt et al., \emph{A Categorical Framework for Learning Generalised Tree Automata}. \url{https://arxiv.org/abs/2001.05786}.
\bibitem{UrbatSchroeder} H. Urbat and L. Schr\"oder, \emph{Automata Learning: An Algebraic Approach}. \url{https://arxiv.org/abs/1911.00874}.
\bibitem{CohenRothblum} J. E. Cohen and U. G. Rothblum, \emph{Nonnegative Ranks, Decompositions, and Factorizations of Nonnegative Matrices}, 1993. \url{https://lab.rockefeller.edu/cohenje/PDFs/211CohenRothblumLAA1993.pdf}.
\bibitem{FritzEtAl} T. Fritz, T. Gonda, P. Perrone and E. F. Rischel, \emph{Representable Markov Categories and Comparison of Statistical Experiments in Categorical Probability}. \url{https://arxiv.org/abs/2010.07416}; \url{https://doi.org/10.1016/j.tcs.2023.113896}.
\end{thebibliography}
\end{document}
